\section{Explicit old-space Euler corrections and the grid cost (NS-60)}
\label{sec:ns60-euler-corrections}

\paragraph{Status and scope.}
The preceding raw-norm obstruction does not apply automatically after
optimal projection. Here we test one explicit family of old-space
corrections. We derive its exact grid error and prove that its full
coefficient norm still cannot obey a fixed-logarithmic $N^{-2}$ bound,
even with suitably growing Euler products. This closes only that uniform
comparison for this correction family. It does not close the optimal
projected comparison, tapered non-Euler mollifiers, or a bound for the
actual residual direction. All results below are analytic; the separate
finite scan is a numerical illustration, not proof evidence.

Extend the difference notation by $\delta_1=\rho_1$. For a finitely
supported integer step function $F$, with value $F_n$ on $(n-1,n]$, set
\begin{equation}
 \Psi F=\sum_{n\geq1}nF_n\delta_n.
 \label{eq:ns60-step-map}
\end{equation}
The triangular change of basis gives
$\operatorname{span}\{\delta_1,\ldots,\delta_N\}=V_N$.
The Mellin identity in Proposition~\ref{prop:v158-mellin-budget} extends
to the first cell because $\int_0^1u^{-s}\,du=(1-s)^{-1}$.

For an integer $d\geq2$, define the old-grid averaging operator
\begin{equation}
 (J_dF)(u)=\int_{m-1}^{m}F(dv)\,dv
 =\frac1d\sum_{d(m-1)<n\leq dm}F_n
 \quad(m-1<u\leq m).
 \label{eq:ns60-cell-average}
\end{equation}
If $F$ is supported in $(N,2N]$, then $J_dF$ is supported in $(0,N]$.
For a squarefree positive integer $Q$ put
\begin{equation}
 \Phi_QF=F+\sum_{\substack{d\mid Q\\d>1}}\mu(d)J_dF,
 \qquad
 \mathcal C_{N,Q}=\text{Gram matrix of }a\longmapsto\Psi\Phi_QF_a,
 \label{eq:ns60-euler-correction}
\end{equation}
where $F_a$ is the new-block step function in
Proposition~\ref{prop:v158-mellin-budget}. When $Q=1$, this is the raw
difference Gram $\mathcal D_N$.

\begin{proposition}[A valid correction, with the grid error retained]
\label{prop:ns60-grid-correction}
The correction lies in $V_N$ and has the explicit old difference
coefficients
\begin{equation}
 (B_{N,Q}a)_m
 =m\sum_{\substack{d\mid Q\\d>1}}\frac{\mu(d)}d
       \sum_{\substack{N<n\leq2N\\d(m-1)<n\leq dm}}\frac{a_n}{n}.
 \label{eq:ns60-correction-matrix}
\end{equation}
Therefore
\begin{equation}
 0<K_N\leq\mathcal C_{N,Q},\qquad
 \mathcal C_{N,Q}
 =\mathcal D_N+X^*B_{N,Q}+B_{N,Q}^*X+B_{N,Q}^*G_\delta B_{N,Q},
 \label{eq:ns60-full-correction-gram}
\end{equation}
where $G_\delta$ is the old difference Gram and
$X=(\langle\delta_m,\delta_n\rangle)_{m\leq N<n\leq2N}$.
Every old/new cross term is present.
Equivalently,
\[
 \mathcal C_{N,Q}-K_N
 =(B_{N,Q}+G_\delta^{-1}X)^*G_\delta
                         (B_{N,Q}+G_\delta^{-1}X).
\]

If $F$ is constant on each interval $(jQ,(j+1)Q]$ on the positive
axis, then $J_dF=F(d\,\cdot)$ for $d\mid Q$. On this compatible
subspace,
\begin{equation}
 \norm{\Psi\Phi_QF}_H^2
 =\frac1{2\pi}\int_{\mathbb R}|\zeta(\tfrac12+it)|^2
       \prod_{p\mid Q}|1-p^{-1/2-it}|^2
       |\widehat q_F(t)|^2\,dt,
 \qquad q_F(v)=e^{v/2}F(e^v).
 \label{eq:ns60-compatible-euler}
\end{equation}
For general $F$, the missing term in this simplification is exactly
\begin{equation}
 R_QF=\sum_{\substack{d\mid Q\\d>1}}\mu(d)
                   \{J_dF-F(d\,\cdot)\}.
 \label{eq:ns60-grid-defect}
\end{equation}
Its Mellin transform and its cross term cannot be omitted.
\end{proposition}
\begin{proof}
Equation~\eqref{eq:ns60-cell-average} gives
\eqref{eq:ns60-correction-matrix}, and the old support gives membership
in $V_N$. Applying $I-P_N$ to $\Psi\Phi_QF_a$ therefore gives the
same projected new vector as applying it to $\Psi F_a$.
Contractivity proves $K_N\leq\mathcal C_{N,Q}$; expanding the squared
norm gives the complete Gram formula.

For compatible $F$, every breakpoint of $F(d\,\cdot)$ is an integer
because $d\mid Q$. Averaging leaves this function unchanged. Also
\[
 \int_0^\infty F(du)u^{-s}\,du
       =d^{s-1}\int_0^\infty F(u)u^{-s}\,du,
 \qquad
 \sum_{d\mid Q}\mu(d)d^{s-1}=\prod_{p\mid Q}(1-p^{s-1}).
\]
On $\Re s=1/2$, the last product has the modulus of
$\prod_{p\mid Q}(1-p^{-s})$. Mellin Plancherel gives
\eqref{eq:ns60-compatible-euler}. This is a finite product identity;
no convergence of an Euler product in the critical strip is assumed.
For general $F$ the transformed
quantity is instead
\[
 \prod_{p\mid Q}(1-p^{s-1})\widehat q_F(t)
                           +\widehat q_{R_QF}(t).
\]
Its absolute square contains both the defect energy and the mixed term.
All these finite step functions have finite weighted Mellin energy,
by the fractional Sobolev argument already proved, so the displayed
identities are Hilbert-space identities rather than formal expansions.
\end{proof}

The grid defect is already nonzero when $N=2$, $Q=2$ and
$F=\mathbf1_{(2,3]}$: then
$F(2u)=\mathbf1_{(1,3/2]}(u)$ whereas
$J_2F=\tfrac12\mathbf1_{(1,2]}$.
Their ordinary $L^2(du)$ distance squared is $1/4$.
This elementary example concerns the grid functions, not a negative
Weil direction or a convergence obstruction.

\begin{lemma}[The smallest Euler multiplier for a bounded grid product]
\label{lem:ns60-euler-minimum}
Let $m_Q=\prod_{p\mid Q}(1-p^{-1/2})$, with $m_1=1$.
For $L=\max(3,\log Q)$,
\begin{equation}
 0\leq-\log m_Q\leq C\frac{\sqrt L}{\log L}
 \label{eq:ns60-euler-minimum}
\end{equation}
with an absolute constant. Consequently, for $Q\leq N$ and sufficiently
large $N$,
$m_Q^{-2}\leq\exp(C\sqrt{\log N}/\log\log N)$.
\end{lemma}
\begin{proof}
For primes $p\geq2$,
$-\log(1-p^{-1/2})\leq C_0p^{-1/2}$.
The elementary prime-counting upper bound
$\pi(x)\ll x/\log x$ (for example
\cite[Corollary 1, Eq. (3.6)]{RosserSchoenfeld1962}) and partial
summation give $\sum_{p\leq L}p^{-1/2}\ll\sqrt L/\log L$.
For primes dividing $Q$ that exceed $L$, use
\[
 \sum_{\substack{p\mid Q\\p>L}}p^{-1/2}
 \leq\frac1{\sqrt L\log L}\sum_{p\mid Q}\log p
 \leq\frac{\sqrt L}{\log L}.
\]
The last inequality also holds when $\log Q<3$. Increasing the absolute
constant covers bounded $L$ and proves the result.
\end{proof}

\begin{proposition}[Growing Euler-grid corrections do not give a logarithmic norm]
\label{prop:ns60-growing-euler-obstruction}
Fix $0\leq\alpha<1$. For each dyadic $N$, let $Q_N$ be any squarefree
positive integer with $Q_N\leq N^\alpha$. For every fixed $B\geq0$,
\begin{equation}
 \norm{\mathcal C_{N,Q_N}}_{\rm op}
 \ne O\left(N^{-2}(\log N)^B\right)
 \quad\text{as }N\to\infty\text{ through powers of two}.
 \label{eq:ns60-growing-no-log}
\end{equation}
The excluded bound is uniform over all new coefficient vectors for the
explicit correction~\eqref{eq:ns60-euler-correction}. It is not a lower
bound for $K_N$, and it does not exclude such a bound restricted to the
actual residual vector $g$.
\end{proposition}
\begin{proof}
Use the fixed nonnegative smooth $w$ from the proof of
Proposition~\ref{prop:v158-lindelof-budget}. Approximate
\[
 F(u)=u^{-1/2}w(\log(u/N))e^{i\tau\log(u/N)}
\]
by its midpoint values on the complete $Q$-cells contained in $(N,2N]$,
and set the other cells to zero. Call the result $F_{N,Q,\tau}$ and set
$a_n=n(F_{N,Q,\tau})_n$. Because $w$ has compact support away from the
two endpoints, for sufficiently large $N$ and $Q\leq N^\alpha$ no
nonzero part is lost in a partially intersecting endpoint cell.
The mean value theorem gives, uniformly in this range,
\begin{equation}
 \norm{q_{F_{N,Q,\tau}}
       -w(\,\cdot-\log N)e^{i\tau(\,\cdot-\log N)}}_{L^1}
 \leq C_w\frac{Q(1+|\tau|)}N,
 \qquad \norm{a}_2^2\leq C_wN^2.
 \label{eq:ns60-packet-grid-error}
\end{equation}
Thus whenever $Q(1+|\tau|)/N$ tends uniformly to zero, the Fourier
transform is bounded below by a fixed positive constant on
$[\tau-1,\tau+1]$. On this compatible packet,
\eqref{eq:ns60-compatible-euler} and
$\prod_{p\mid Q}|1-p^{-1/2-it}|\geq m_Q$ imply
\begin{equation}
 \int_{\tau-1}^{\tau+1}|\zeta(\tfrac12+it)|^2\,dt
 \leq C m_Q^{-2} N^2\norm{\mathcal C_{N,Q}}_{\rm op}.
 \label{eq:ns60-corrected-local-test}
\end{equation}

Suppose the excluded big-$O$ bound held. For large $T$, choose a power
of two $N$ with
$T^{2/(1-\alpha)}\leq N<2T^{2/(1-\alpha)}$.
Then $Q_N(1+2T)/N\leq(1+2T)/T^2\to0$, uniformly for
$|\tau|\leq2T$. Cover the interval in
\eqref{eq:v158-short-interval} by $O(\log^2T)$ unit intervals and use
\eqref{eq:ns60-corrected-local-test}. Lemma~\ref{lem:ns60-euler-minimum}
gives
\[
 |\zeta(\tfrac12+iT)|^2
 \ll \log T+
 \exp\left(C_\alpha\frac{\sqrt{\log T}}{\log\log T}\right)
 (\log T)^{B+3}.
\]
The logarithm of the right side is
$O_\alpha(\sqrt{\log T}/\log\log T)+(B+3)\log\log T$.
It is $o(\sqrt{\log T/\log\log T})$, contradicting the unconditional
large values in~\cite[Theorem 1]{Soundararajan2008}.
No relation between the choices $Q_N$ at different $N$ is required.
\end{proof}

\paragraph{What changed and what remains open.}
Downsampling does give an admissible old-space correction, but on a
general new vector it requires a cell average. Replacing that average
by an exact Euler multiplier loses~\eqref{eq:ns60-grid-defect} and its
mixed term. Even on a compatible family where the replacement is exact,
known zeta large values defeat the stated uniform norm goal whenever
$Q_N\leq N^\alpha$ for fixed $\alpha<1$.
Since $K_N\leq\mathcal C_{N,Q_N}$, this failure of the explicit upper
comparison does not give a lower obstruction for the optimal projection.
The range $Q_N=N^{1-o(1)}$, tapered non-Euler coefficients and
actual-residual-only estimates are outside the proved exclusion.

The separately labelled scan in
\texttt{evidence/diag\_ns60\_nb\_corrections/} compares fixed Euler
products $Q=2,6,30$ and tapered M\"obius lengths $4,16,64$ with the full
projection on fixed blocks $32\leq N\leq512$. It tests the correction
proposal only; its finite reductions and selected-direction values are
not asymptotic bounds. The remaining RH-relevant input is still a
cofinal gain estimate for the actual optimal residuals. No such estimate,
NB convergence, G2 or RH proof is supplied.

\section{One Mellin integration: an elementary block budget (NS-61)}
\label{sec:ns61-smoothed-nb}

\paragraph{Status and prior art.}
The $q=2$ weighted Nyman--Beurling problem already appears in
\cite{Ehm2024}; it is not a new RH criterion. We prove its equivalence
here with the tail space stated, and use adjacent differences to obtain
an unconditional $N^{-2}$ block budget. The numerator is transformed
along with the norm. No asymptotic lower correlation estimate, NB
convergence, G2 or RH proof follows from this upper budget.

Let $H=L^2(0,\infty;dx)$ and retain $\rho_n$, $\chi$, $V_N$ and
$\kappa=\log(2\pi)-\gamma$ from Section~\ref{sec:ns53-nb-blocks}.
Define the bounded Mellin convolution
\[
 (Jf)(x)=\int_0^1f(x/u)\,\frac{du}{u},\qquad
 \sigma_n=J\rho_n,\qquad \tau=J\chi=(-\log x)\mathbf1_{(0,1)}.
\]
On the Mellin line $\Re s=1/2$, $J$ has multiplier $1/s$, so
$\norm J=2$ and $J$ is injective. Set
\[
 W_N=\operatorname{span}\{\sigma_1,\ldots,\sigma_N\},\quad
 R_N=\tau-\Pi_N\tau,\quad E_N^{(2)}=\norm{R_N}^2,
\]
where $\Pi_N$ is the orthogonal projection onto $W_N$.
The superscript is a smoothing label; $E_N^{(2)}$ is already a
squared norm, and $\norm\tau^2=2$.

\begin{proposition}[The smoothed target still detects an off-line zero]
\label{prop:ns61-rh-equivalence}
Every finite $E_N^{(2)}$ is positive, and
\[
 E_N^{(2)}\longrightarrow0\quad\Longleftrightarrow\quad\mathrm{RH}.
\]
If $\zeta(\beta+i\gamma_*)=0$ with $1/2<\beta<1$, put
$z=\beta+i\gamma_*$. Then, for every $N$,
\begin{equation}
 E_N^{(2)}\geq
 \frac{|z|^{-4}}{(2\beta-1)^{-1}+|1-z|^{-2}}>0.
 \label{eq:ns61-zero-obstruction}
\end{equation}
\end{proposition}
\begin{proof}
The original strong NB criterion and boundedness of $J$ give
$E_N^{(2)}\leq4E_N$, proving the forward approximation under RH.
For the converse, the functions $\tau$ and $\sigma_n$ belong to the
closed space
\[
 \mathcal W=L^2(0,1)\oplus
       \operatorname{span}\{x^{-1}\mathbf1_{(1,\infty)}\}.
\]
Indeed $\sigma_n(x)=1/(nx)$ for $x>1$. On this space the functional
\[
 L_z(f)=\int_0^1 f(x)x^{z-1}\,dx+\frac{c}{1-z},
 \qquad f(x)=c/x\quad(x>1),
\]
is bounded with norm squared $(2\beta-1)^{-1}+|1-z|^{-2}$.
The absolutely convergent Mellin identities on $0<\Re s<1$ are
\[
 \mathcal M\sigma_n(s)=-\frac{n^{-s}\zeta(s)}{s^2},
 \qquad \mathcal M\tau(s)=\frac1{s^2}.
\]
Thus $L_z(\sigma_n)=0$ and $L_z(\tau)=z^{-2}$, which proves
\eqref{eq:ns61-zero-obstruction}. If RH fails, the functional equation
supplies a zero on the right of the critical line. This proves the
converse without assuming a bounded inverse for $J$.
Injectivity of $J$ also transfers finite linear independence and
$\chi\notin V_N$ from the original finite problem, so $E_N^{(2)}>0$.
\end{proof}

\begin{proposition}[A unitary average identity and an elementary trace budget]
\label{prop:ns61-average-budget}
Put $U=J-I$ and $h=U\rho_1$. The operator $U$ is unitary on $H$,
$\norm h^2=\kappa$, and $h$ is nonnegative and supported in $(0,1]$.
For $n\geq2$ define
\[
 \epsilon_n=\sigma_n-\frac{n-1}{n}\sigma_{n-1}=J\delta_n.
\]
Then, as a Bochner integral in $H$,
\begin{equation}
 n\epsilon_n=\int_{n-1}^{n}h(u\,\cdot)\,du
       =U\int_{n-1}^{n}\rho_1(u\,\cdot)\,du,
 \qquad
 \norm{\epsilon_n}^2\leq
       \frac{\kappa}{n^2}\log\frac n{n-1}.
 \label{eq:ns61-average}
\end{equation}
In particular $\epsilon_n$ is nonnegative and supported in
$(0,1/(n-1)]$, with
\begin{equation}
 \int_0^\infty x\epsilon_n(x)\,dx
             =\frac{\pi^2}{24n^2(n-1)}.
 \label{eq:ns61-positive-first-moment}
\end{equation}
If $\mathcal D_N^{(2)}$ is the raw Gram of
$\epsilon_{N+1},\ldots,\epsilon_{2N}$, then
\begin{equation}
 \norm{\mathcal D_N^{(2)}}_{\rm op}
 \leq\operatorname{tr}\mathcal D_N^{(2)}\leq \mathfrak b_N,
 \qquad \mathfrak b_N=\kappa\sum_{n=N+1}^{2N}
          \frac{\log(n/(n-1))}{n^2}
 \leq\frac{3\kappa}{8N^2}.
 \label{eq:ns61-trace-budget}
\end{equation}
Moreover $N^2\operatorname{tr}\mathcal D_N^{(2)}\to3\kappa/8$.
These are unconditional analytic bounds and a trace asymptotic, not
an asymptotic formula for the residual gain.
\end{proposition}
\begin{proof}
The Mellin multiplier of $U$ is $(1-s)/s$, which has modulus one on
the critical line and is nonzero there. Hence $U$ is unitary.
Writing
\[
 A(z)=\int_0^z\frac{\{v\}}v\,dv,
 \qquad \sigma_u(x)=A(1/(ux))\quad(u>0),
\]
gives $u\partial_u\sigma_u=-\rho_u$ almost everywhere. Therefore
$\partial_u(u\sigma_u)=\sigma_u-\rho_u=h(u\,\cdot)$.
On compact positive $u$ intervals this is also the strong integral
identity: dilations are strongly continuous in $H$, and integrating
the displayed derivative recovers the pointwise formula and an
$H$-integrable majorant. Since $\sigma_1(x)=\rho_1(x)=1/x$ for
$x>1$, $h$ has the asserted support.
For $z\in[k,k+1)$ with $k\geq1$, integration in the definition of
$A$ gives $A(z)=z-k\log z+\log(k!)$, hence
$h(1/z)=k-k\log z+\log(k!)>0$.
Indeed monotonicity of $\log u$ gives
$\int_1^{k+1}\log u\,du\leq\log((k+1)!)$, or
$\log(k!)\geq k\log(k+1)-k$; use $z<k+1$.
Positivity of $\epsilon_n$ follows from its average representation.
The compactly supported $h$ has Mellin transform
$-(1-s)\zeta(s)/s^2$, which extends analytically to $\Re s>0$.
At $s=2$ it gives $\int xh(x)\,dx=\pi^2/24$. Integrating the
average identity against $x\,dx$ proves
\eqref{eq:ns61-positive-first-moment}.
The Cauchy--Schwarz inequality for a unit-length Bochner integral gives
\[
 n^2\norm{\epsilon_n}^2\leq
 \int_{n-1}^{n}\norm{h(u\,\cdot)}^2\,du
 =\kappa\int_{n-1}^{n}\frac{du}{u}.
\]
For $n-1\leq u\leq n$, $1/(n^2u)\leq u^{-3}$, whence
\[
 \mathfrak b_N\leq\kappa\int_N^{2N}u^{-3}\,du
       =\frac{3\kappa}{8N^2}.
\]
Finally, dilation continuity gives
$n^3\norm{\epsilon_n}^2\to\kappa$: rescale $x=y/n$ in
\eqref{eq:ns61-average}, then $h((u/n)y)\to h(y)$ strongly,
uniformly for $n-1\leq u\leq n$. Summation on $N<n\leq2N$
and $N^2\sum n^{-3}\to3/8$ prove the trace asymptotic.
\end{proof}

The associated exact Gram identity, which retains all off-diagonal
entries, is
\begin{equation}
 \langle\epsilon_m,\epsilon_n\rangle
 =\frac1{mn}\int_{m-1}^{m}\int_{n-1}^{n}
        \langle\rho_u,\rho_v\rangle\,du\,dv.
 \label{eq:ns61-averaged-full-gram}
\end{equation}
Thus smoothing converts the singular adjacent-difference budget into
an average of the original complete Gram kernel. It does not diagonalize
that kernel or remove the old-space projection.

\begin{proposition}[Complete smoothed gain and its still-missing numerator]
\label{prop:ns61-gain}
Form the old Gram $G^{(2)}$, old/new block $C^{(2)}$, new Gram
$D^{(2)}$, and target loads $b^{(2)}$ using $\sigma_n$ and $\tau$.
Set
\[
 c=(G^{(2)})^{-1}b^{(2)},\quad
 h^{(2)}=b_{\rm new}^{(2)}-(C^{(2)})^*c,\quad
 S^{(2)}=D^{(2)}-(C^{(2)})^*(G^{(2)})^{-1}C^{(2)}.
\]
For the same triangular new-coordinate matrix $T_N$ as in
Proposition~\ref{prop:v158-difference-gain}, put
$g^{(2)}=T_N^*h^{(2)}$ and $K_N^{(2)}=T_N^*S^{(2)}T_N$.
Then
\begin{equation}
 0<K_N^{(2)}\leq\mathcal D_N^{(2)},\qquad
 E_N^{(2)}-E_{2N}^{(2)}
 =(g^{(2)})^*(K_N^{(2)})^{-1}g^{(2)}
 \geq\frac{\norm{g^{(2)}}_2^2}{\mathfrak b_N}
 \geq\frac{8N^2}{3\kappa}\norm{g^{(2)}}_2^2.
 \label{eq:ns61-gain}
\end{equation}
Consequently the cofinal estimate
\begin{equation}
 N^2\norm{g_{2^j}^{(2)}}_2^2
 \geq\frac{3\kappa}{8}\alpha_j E_{2^j}^{(2)},
 \quad N=2^j,\quad 0\leq\alpha_j<1,\quad
 \sum_j\alpha_j=\infty
 \label{eq:ns61-open-correlation}
\end{equation}
for all sufficiently large $j$ would prove RH. This lower estimate
is an open input; no necessity claim is made for it.
\end{proposition}
\begin{proof}
The finite Gram matrices are positive definite by injectivity of $J$.
Orthogonal projection, the complete Schur complement and the invertible
coordinate change give the exact identities as in NS-53.
The raw trace budget bounds the largest eigenvalue of the projected
Gram. The resulting contraction
$E_{2^{j+1}}^{(2)}\leq(1-\alpha_j)E_{2^j}^{(2)}$ proves the
conditional conclusion. In particular, the unsmoothed residual
correlation cannot be inserted in~\eqref{eq:ns61-gain}.
\end{proof}

\begin{proposition}[Explicit loads do not supply a projected lower bound]
\label{prop:ns61-loads}
Use the standard convention
$\zeta(1+z)=z^{-1}+\gamma-\gamma_1z+O(z^2)$ for $\gamma_1$.
Then
\begin{equation}
 b_n^{(2)}=\langle\tau,\sigma_n\rangle
 =\frac{P(\log n)}n,
 \quad P(v)=\tfrac12v^2+(2-\gamma)v+3-2\gamma-\gamma_1.
 \label{eq:ns61-loads}
\end{equation}
In difference coordinates, for $n\geq2$,
\[
 \langle\tau,\epsilon_n\rangle
 =\frac{P(\log n)-P(\log(n-1))}{n}
 \sim\frac{\log n+2-\gamma}{n^2}.
\]
For the actual residual the exact expression is instead
\begin{equation}
 g_n^{(2)}=\frac{P(\log n)-P(\log(n-1))}{n}
       -\sum_{m\leq N}c_m\langle\sigma_m,\epsilon_n\rangle,
 \quad N<n\leq2N.
 \label{eq:ns61-retained-cross}
\end{equation}
The sign and size of the second term are not controlled here.
\end{proposition}
\begin{proof}
Since $\sigma_1(x)=1/x$ for $x>1$, the transform
$-\zeta(s)/s^2-1/(1-s)$ is the Mellin transform of the restriction
to $(0,1)$. Its value and derivative at $s=1$ give
\[
 \int_0^1\sigma_1(x)\,dx=2-\gamma,\qquad
 \int_0^1(-\log x)\sigma_1(x)\,dx=3-2\gamma-\gamma_1.
\]
Change variables $u=nx$ in the load integral and split at $u=1$.
The tail contribution is $(\log n)^2/(2n)$, giving
\eqref{eq:ns61-loads}. Taking the adjacent difference and then
subtracting the old projection gives the remaining identities.
The loads agree, with the sign convention for $\sigma_n$ reversed,
with the $q=2$ formula in~\cite{Ehm2024}.
\end{proof}

\paragraph{The exact point at which the lower-bound attempt stops.}
Every new $\epsilon_n$ is supported in $(0,1/N]$, so the numerator
only sees the small-$x$ part of $R_N$. The elementary formula
\[
 A(z)=z-k\log z+\log(k!),\qquad k=\lfloor z\rfloor,
\]
and Stirling's estimate imply, uniformly for $z\geq1$,
$|A(z)-\tfrac12\log(2\pi z)|\leq3/z$; a complete elementary
remainder is given in Section~\ref{sec:ns61-gram-remainder}.
Thus for $0<x\leq1/N$,
\begin{equation}
 \begin{aligned}
 R_N(x)&=a_N\log(1/x)+\ell_N+\mathcal R_N(x),\\
 a_N&=1-\tfrac12\sum_{m\leq N}c_m,\qquad
 \ell_N=\tfrac12\sum_{m\leq N}c_m\log\frac m{2\pi},\\
 |\mathcal R_N(x)|&\leq3x\sum_{m\leq N}m|c_m|.
 \end{aligned}
 \label{eq:ns61-residual-endpoint}
\end{equation}
Writing $\Delta_n f=f(n)-f(n-1)$, the exact first two moments of
$\epsilon_n$ give
\begin{equation}
 g_n^{(2)}=
 \frac{a_N\{\tfrac12\Delta_n(\log^2)+(2-\gamma)\Delta_n\log\}
              +\ell_N\Delta_n\log}{n}
 +\mathcal E_{N,n},
 \quad N<n\leq2N,
 \label{eq:ns61-endpoint-correlation}
\end{equation}
where
\[
 |\mathcal E_{N,n}|\leq
       \frac{\pi^2\sum_{m\leq N}m|c_m|}{8n^2(n-1)}.
\]
The two leading moments follow from $\int h=1$ and
$\int(-\log x)h(x)\,dx=2-\gamma$, obtained by taking the value
and derivative of $-(1-s)\zeta(s)/s^2$ at $s=1$.
For the error, use nonnegativity and the exact first moment
\eqref{eq:ns61-positive-first-moment}. Neither a lower bound for the displayed
coefficient combination relative to $E_N^{(2)}$, nor a small enough
bound for $\sum m|c_m|$, has been obtained. Replacing these optimized
coefficients by the target loads would discard
\eqref{eq:ns61-retained-cross}.

The improvement is therefore a proved elementary denominator and an
explicit description of the remaining arithmetic numerator. It bypasses
the raw unsmoothed norm obstruction by changing both the functions and
the target. It does not bypass the zero obstruction, certify a lower
correlation, or establish convergence. This is the stop condition for
NS-61 unless an independent estimate for the optimized residual is found.

\subsection{A complete remainder for the smoothed Gram computation}
\label{sec:ns61-gram-remainder}

This subsection supplies the analytic tail used by the finite Arb replay;
the diagnostic truncations do not enter that replay. Write
$\widetilde B_j(x)=B_j(\{x\})$ and let $R_1,R_2,S_2$ be the
functions in~\cite[Theorem 2.1]{Ehm2024}. For $r=p/q\geq1$ in lowest
terms the Gram formula used is
\[
 G^{(2)}_{m,n}=\frac{K_2+K_1\log(n/m)+\tfrac14\log^2(n/m)}n
                  +\frac{S_2(n/m)}m\quad(m\leq n),
\]
with
\[
 K_1=\frac{\log(2\pi)-\gamma+2}{2},\qquad
 K_2=(1-\gamma/2)\log(2\pi)+\tfrac14\log^2(2\pi)
       +\frac{\pi^2}{48}-\frac{\gamma^2}{4}-\gamma-\gamma_1+\frac32.
\]
The first $K$ terms of $S_2(r)=\sum_{k\geq1}R_2(kr)$ are evaluated
by the finite formula in that theorem, retaining balls throughout.
The rest is enclosed by the following lemma.

\begin{lemma}[Bernoulli tail with an explicit remainder]
\label{lem:ns61-r2-tail}
For an integer $M\geq4$, let $c_2=c_3=0$ and, for $j\geq4$, let
\[
 c_j=\frac{2j-3-(j-1)H_{j-1}}{j(j-1)},\qquad
 P_M(x)=\sum_{j=4}^{M+2}c_j\frac{\widetilde B_j(x)}{x^j}.
\]
Put
\[
 a_{M+1}=2(M+2)(M+1)c_{M+2}-(M+1)Mc_{M+1},\qquad
 a_{M+2}=-(M+2)(M+1)c_{M+2}.
\]
For any constants $U_j\geq\sup_{0\leq x\leq1}|B_j(x)|$,
\begin{equation}
 |R_2(x)-P_M(x)|\leq
 \frac{U_M}{M^2(M-1)x^M}
 +\frac{|a_{M+1}|U_{M+1}}{M(M+1)x^{M+1}}
 +\frac{|a_{M+2}|U_{M+2}}{(M+1)(M+2)x^{M+2}}.
 \label{eq:ns61-r2-remainder}
\end{equation}
The rational constants
$U_j=2j!\,(1+1/(j-1))/6^j$, $j\geq2$, are valid.
For $r=p/q$ the Bernoulli part of the infinite sum is exactly
\begin{equation}
 \sum_{k>K}\frac{\widetilde B_j(kr)}{(kr)^j}
 =p^{-j}\sum_{a=1}^{q}
 B_j\left(\left\{\frac{(K+a)p}{q}\right\}\right)
 \zeta\left(j,\frac{K+a}{q}\right).
 \label{eq:ns61-hurwitz-tail}
\end{equation}
The sum of the three remainders in~\eqref{eq:ns61-r2-remainder}
is bounded by replacing each $x^{-l}$ by
$r^{-l}\zeta(l,K+1)$.
\end{lemma}
\begin{proof}
From $R_1(x)=-\int_x^\infty\widetilde B_1(t)t^{-2}\,dt$,
successive integrations by parts give the exact finite formula
\[
 R_1(x)=\sum_{j=2}^{M}\frac{\widetilde B_j(x)}{jx^j}
            -\int_x^\infty\frac{\widetilde B_M(t)}{t^{M+1}}\,dt.
\]
The remainder is at most $U_M/(Mx^M)$. Also
\cite[Proposition 5.2, Eq. (40)]{Ehm2024} gives
$-(x^2R_2'(x))'=R_1(x)$ and the inverse formula
\[
 (\mathcal If)(x)=\int_x^\infty f(t)\left(\frac1t-\frac1x\right)dt.
\]
If $|f(t)|\leq At^{-l}$ with $l>1$, then
$|\mathcal If(x)|\leq A x^{-l}/(l(l-1))$.
For $\mathcal L=-(x^2\partial_x)'$ the coefficient of
$\widetilde B_k(x)/x^k$ in $\mathcal LP_M$ is
\[
 -(k+2)(k+1)c_{k+2}+2(k+1)kc_{k+1}-k(k-1)c_k.
\]
It equals $1/k$ for $2\leq k\leq M$; the two remaining coefficients
are $a_{M+1},a_{M+2}$. Periodic Bernoulli polynomials of orders at
least four are twice continuously differentiable across integers, so
no point masses are omitted in this differentiation. Apply
$\mathcal I$ to $R_1-\mathcal LP_M$ and the preceding three bounds.
The possible homogeneous terms $A+B/x$ vanish because both $R_2$
and $P_M$ are $O(x^{-4})$. This proves~\eqref{eq:ns61-r2-remainder}.

The absolutely convergent Bernoulli Fourier series
\cite[Eq. (24.8.3)]{DLMFBernoulliFourier} gives
$\sup|B_j|\leq2j!\zeta(j)/(2\pi)^j$.
Using $2\pi>6$ and $\zeta(j)\leq1+1/(j-1)$ proves the rational
upper bound. Finally write $k=K+a+q\ell$ and sum over $\ell\geq0$.
The periodic numerator is constant in each such progression, giving
\eqref{eq:ns61-hurwitz-tail}. Absolute convergence justifies every
rearrangement and the sum of the absolute remainder bounds.
\end{proof}

\paragraph{An independent positive-cell convention check.}
On a cell where $k=\lfloor t/m\rfloor$,
$A(t/m)=t/m-k\log t+k\log m+\log(k!)$.
Products of two such expressions divided by $t^2$ are integrated
exactly between consecutive multiples of $m$ and $n$; the initial
cell has constant integrand $1/(mn)$.
For $z=k+r\geq1$, $0\leq r<1$, the elementary Stirling remainder
\cite[\S5.11]{DLMFGammaBounds}
$0<\log(k!)-(k+1/2)\log k+k-\tfrac12\log(2\pi)<1/(12k)$
gives
\[
 \left|A(z)-\tfrac12\log(2\pi z)\right|
 \leq\frac{13}{12k}\leq\frac{13}{6z}<\frac3z.
\]
Indeed $0\leq r-k\log(1+r/k)\leq r^2/(2k)$ and
$0\leq\tfrac12\log(1+r/k)\leq r/(2k)$.
For $T\geq\max(m,n)$ put $L_m=\tfrac12\log(2\pi T/m)$.
The remaining physical integral is within
\[
 \frac3{T^2}\left[m\left(\frac{L_n}2+\frac18\right)
                 +n\left(\frac{L_m}2+\frac18\right)\right]
            +\frac{3mn}{T^3}
\]
of $(L_mL_n+(L_m+L_n)/2+1/2)/T$.
This separate integral enclosure checks normalization and several
off-diagonal entries of the series formula; it is not used to sharpen
the block estimates.

\paragraph{Certified finite checks, not a cofinal estimate.}
The complete smoothed Grams through index 64 were enclosed at 256 and
384 bits with the entire series tail retained. A separate replay doubles
the finite series cutoff from 128 to 256 while keeping the same matrix.
All 180 enclosure comparisons across the precision and cutoff replays pass; the minimum internal accuracies are
99, 99 and 123 bits. Seven independent positive-cell Gram integrals
also pass at both precisions.
\begin{center}
\begin{tabular}{rccc}
\toprule
$N$ & exact relative gain & trace-budget relative bound & bound/gain\\
\midrule
16 & $0.2372848\ldots$ & $0.003293361\ldots$ & $0.01387935\ldots$\\
32 & $0.1891280\ldots$ & $0.003566353\ldots$ & $0.01885681\ldots$\\
\bottomrule
\end{tabular}
\end{center}
The trace lower estimate is below $1/50$ of the exact gain on both
blocks. The endpoint remainder budget is more than 2000 times the norm
of its leading term there, so that particular triangle lower bound is
vacuous, despite every individual remainder gate passing.
These are finite statements, not a failure of
\eqref{eq:ns61-open-correlation} or an asymptotic claim.

\section{Smoothing does not repair the canonical sharp interpolant (NS-62)}
\label{sec:ns62-sharp-interpolant}

\paragraph{Status and scope.}
The previous section changes the approximation norm and proves a smaller
upper budget. It is natural to ask whether simply smoothing the canonical
M\"obius interpolant now supplies convergence. It does not. The following
localized dual test extends the classical sharp-interpolant obstruction
in~\cite[Section 4]{BaezDuarte2000Arithmetic} through Mellin integration.
The conclusion concerns this explicit coefficient rule, including its
tail correction. It does not concern the optimized coefficients in
Proposition~\ref{prop:ns61-gain}.

Define
\[
 M(x)=\sum_{k\leq x}\mu(k),\quad
 s_N=\sum_{k\leq N}\frac{\mu(k)}k,\quad
 \eta_N=s_N-\frac{M(N)}N
       =\sum_{k<N}\frac{M(k)}{k(k+1)},
\]
and let
\begin{equation}
 B_N=\sum_{k=1}^N\mu(k)\rho_k-Ns_N\rho_N.
 \label{eq:ns62-natural-interpolant}
\end{equation}
The sign is chosen to match the classical notation: $B_N=-\chi$
on $x>1/N$ (up to endpoints). The sum is zero on $x>1$.
The smoothed target error is $J^r(\chi+B_N)$.

\begin{proposition}[A localized witness survives any number of integrations]
\label{prop:ns62-smoothed-obstruction}
For every integer $N\geq1$ and integer $r\geq0$,
\begin{equation}
 \norm{J^r(\chi+B_N)}_H
 \geq\frac{|N\eta_N-k_{N,r}|}{\sqrt N}
 \geq\frac{|N\eta_N-2|}{\sqrt N}
            -\frac{4\pi^2}{3^{r+2}N^{5/2}},
 \label{eq:ns62-laguerre-lower}
\end{equation}
where
\[
 k_{N,r}=\frac N{r!}\int_0^{1/N}
       \frac{\sin(2\pi x)}{\pi x}
                    \left(\log\frac1{Nx}\right)^r dx.
\]
In particular, for every sequence of nonnegative integers $r_N$,
$\norm{J^{r_N}(\chi+B_N)}_H$ does not tend to zero.
For growing $r_N$ the target also changes with $N$; this assertion
is only a failure of these particular approximants, not a new RH
criterion for varying norms.
The nonconvergence is for the full sequence of integers $N$. No
exclusion of a selected subsequence, including the dyadic subsequence,
is asserted without a corresponding lower bound for $\sqrt N|\eta_N|$
on that subsequence.
\end{proposition}
\begin{proof}
Use the dilation-invariant unitary $\mathcal V$ whose critical-line
Mellin multiplier, away from the discrete zero set, is
\[
 v(s)=\frac{s\zeta(1-s)}{(1-s)\zeta(s)}.
\]
Its modulus is one on $\Re s=1/2$; define it arbitrarily with modulus
one on the measure-zero exceptional set. Mellin identities and the
functional equation give
\[
 \mathcal V\rho_k(x)=\frac{\{kx\}}{kx},\qquad
 \mathcal V\chi(x)=k(x):=\frac{\sin(2\pi x)}{\pi x}.
\]
These are also the defining identities of the classical unitary in
\cite[Section 4]{BaezDuarte2000Arithmetic}. Both $\mathcal V$ and $J$
are Mellin multipliers, so they commute. For $0<x<1/N$,
\begin{equation}
 \mathcal V(\chi+B_N)(x)=k(x)+M(N)-Ns_N=k(x)-N\eta_N.
 \label{eq:ns62-unitary-small-cell}
\end{equation}

Let $L_r(z)=\sum_{j=0}^r(-1)^j\binom rj z^j/j!$ be the ordinary
Laguerre polynomial and put
\[
 \psi_{N,r}(x)=\sqrt N\,
       L_r\left(\log\frac1{Nx}\right)\mathbf1_{(0,1/N)}(x).
\]
Laguerre orthogonality \cite[Table 18.3.1]{DLMFLaguerre} gives
$\norm{\psi_{N,r}}=1$. The adjoint is
$(J^*f)(x)=x^{-1}\int_0^x f(y)\,dy$. Repeated integration, or
Rodrigues' formula and integration by parts, gives the complete
identity
\begin{equation}
 (J^*)^r\psi_{N,r}(x)=
 \frac{(-1)^r\sqrt N}{r!}
       \left(\log\frac1{Nx}\right)^r\mathbf1_{(0,1/N)}(x).
 \label{eq:ns62-adjoint-laguerre}
\end{equation}
There is no exterior term: at the $j$th integration, $j\leq r$,
the potential exterior coefficient is a Laguerre moment of a
polynomial of degree less than $r$, hence zero. On the interior,
the polynomial identity follows by integrating monomials against
$e^{-z}$; equivalently $(1-\partial_z)^{-r}L_r(z)=(-1)^rz^r/r!$.
This also covers $r=0$ directly.

Taking the inner product with $\psi_{N,r}$, moving $J^r$ to the
adjoint, and using~\eqref{eq:ns62-unitary-small-cell} yields
\[
 \left|\left\langle J^r\mathcal V(\chi+B_N),\psi_{N,r}\right\rangle\right|
       =\frac{|k_{N,r}-N\eta_N|}{\sqrt N}.
\]
The unitary norm and Cauchy--Schwarz prove the first inequality.
For real $x$, $|k(x)-2|\leq4\pi^2x^2/3$. Since
\[
 \int_0^{1/N}x^2\left(\log\frac1{Nx}\right)^r dx
        =\frac{r!}{3^{r+1}N^3},
\]
we have $|k_{N,r}-2|\leq4\pi^2/(3^{r+2}N^2)$, proving the second.

For completeness, $\eta_N\ne o(N^{-1/2})$ follows from the known
existence of a critical-line zeta zero. Indeed let
$\eta(x)=\int_1^xM(t)t^{-2}\,dt$. For $\Re s>0$,
\[
 \int_1^\infty\eta(x)x^{-s-1}\,dx
        =\frac1{s(s+1)\zeta(s+1)}.
\]
The right side has a genuine pole at $s=\rho-1$ for any
critical-line zero $\rho$. If $\eta(x)=o(x^{-1/2})$, the integral
would be analytic on $\Re s>-1/2$, and the usual Abelian boundary
estimate would make $(\sigma+1/2)$ times its value at
$\sigma+i\Im\rho$ tend to zero as $\sigma\downarrow-1/2$.
This contradicts a pole of any positive order at that point.
Also $\eta(N)=\eta_N$ and
$\eta(x)-\eta(\lfloor x\rfloor)=O(1/x)$, so the discrete
little-$o$ assertion would imply the continuous one. Thus
$\limsup_N\sqrt N|\eta_N|>0$, whereas the error term in
\eqref{eq:ns62-laguerre-lower} tends to zero uniformly in $r\geq0$.
This proves the stated nonconvergence for every choice of $r_N$.
No assumption about zeros off the critical line was made.
\end{proof}

\begin{proposition}[Positive smoothed differences still require signed coefficients]
\label{prop:ns62-positive-cone}
Every $\epsilon_n$, $n\geq2$, is nonnegative. Nevertheless $\tau$
is not in the $H$-closure of the cone of finite nonnegative linear
combinations of $\sigma_1,\epsilon_2,\epsilon_3,\ldots$.
This is a restriction on that coefficient cone, not on their signed
linear span or on the Weil form.
\end{proposition}
\begin{proof}
Nonnegativity was proved in Proposition~\ref{prop:ns61-average-budget}.
Its average identity also gives strict positivity on
$(1/n,1/(n-1))$.

Triangular telescoping gives
\begin{equation}
 -JB_N=\sum_{n=2}^{N}n s_{n-1}\epsilon_n.
 \label{eq:ns62-pointwise-series}
\end{equation}
It equals $\tau$ on $x>1/N$, because $-B_N=\chi$ there and
$Jf(x)$ uses only values at arguments at least $x$.
For $N=6$ the coefficient of $\epsilon_6$ is
$6s_5=6(1-1/2-1/3-1/5)=-1/5$.

Restrict now to $x>1/6$. Every $\epsilon_n$ with $n\geq7$
vanishes there. The restrictions of
$\sigma_1,\epsilon_2,\ldots,\epsilon_6$ are linearly independent:
first use $x>1$ for $\sigma_1$, then the successive intervals
$(1/n,1/(n-1))$ for $n=2,\ldots,6$.
Their finite-dimensional nonnegative cone is closed, but the unique
representation of the restricted target has the negative coefficient
$-1/5$. Thus the restricted target has positive distance from this
cone. Restriction is contractive, proving the full-space claim.
\end{proof}

\paragraph{Consequences for the next attempt.}
The positive shapes of the smoothed differences do not make the
coefficient problem positive. The canonical pointwise interpolation
coefficients are explicit, but their partial sums still fail in norm
after smoothing. The optimized residual criterion from NS-61 remains
open and is unaffected by this exclusion. A next attempt must obtain
an independent estimate for those adaptive residuals, or a different
coefficient rule with a complete norm bound; visual smoothness,
pointwise exactness away from zero, and a small raw block norm do not
supply it.

\section{Fixed smoothing and a changing-norm control (NS-63)}
\label{sec:ns63-changing-norm}

\paragraph{Why this check is needed.}
NS-61 uses one fixed extra Mellin integration. Increasing the smoothing
order with the approximation size changes the problem. The following
control shows that relative errors can then vanish with just one fixed
dilation and a fixed coefficient, unconditionally. It is not convergence
in a fixed NB norm.

For an integer $r\geq0$ put
\[
 \tau_r=J^r\chi=\frac{(-\log x)^r}{r!}\mathbf1_{(0,1)},\quad
 \sigma_{n,r}=J^r\rho_n,\quad
 T_r=\norm{\tau_r}^2=\binom{2r}{r}.
\]
Let $E_{N,r}$ be the squared distance from $\tau_r$ to the span of
$\sigma_{1,r},\ldots,\sigma_{N,r}$, so $E_{N,1}=E_N^{(2)}$.

\begin{proposition}[Keep the smoothing order fixed in the RH criterion]
\label{prop:ns63-fixed-order}
For each fixed $r$, $E_{N,r}\to0$ if and only if RH. If
$z=\beta+i\gamma_*$ is a zero with $1/2<\beta<1$, then
\[
 E_{N,r}\geq
 \frac{|z|^{-2(r+1)}}{(2\beta-1)^{-1}+|1-z|^{-2}}.
\]
The lower obstruction itself depends on $r$. After division by $T_r$
it tends to zero for every such $z$, since $|z|>1/2$; it gives no
fixed positive obstruction for the relative error as $r\to\infty$.
\end{proposition}
\begin{proof}
Boundedness of $J^r$ gives $E_{N,r}\leq4^rE_N$ under the original
strong NB criterion. Every $\sigma_{n,r}$ has the same tail $1/(nx)$
for $x>1$, whereas $\tau_r$ has zero tail. Thus the bounded functional
$L_z$ in Proposition~\ref{prop:ns61-rh-equivalence} applies unchanged.
Its value on the target is $z^{-(r+1)}$ and it annihilates the finite
span. This proves the lower bound and the converse for fixed $r$.
\end{proof}

\begin{proposition}[One fixed atom has vanishing relative error as the norm changes]
\label{prop:ns63-one-atom}
Let $z_0=\zeta(1/2)$, $c_0=-1/z_0$, and
\[
 D_0=\frac{\zeta'(1/2)}{\zeta(1/2)}
     =\frac12\left(\gamma+\frac\pi2+\log(8\pi)\right)>0.
\]
Unconditionally, as $r\to\infty$ through integers,
\begin{equation}
 \frac{\norm{\tau_r-c_0\sigma_{1,r}}^2}{T_r}
    =\frac{D_0^2}{4(2r-1)}+O(r^{-2})
    \sim\frac{D_0^2}{8r}\longrightarrow0.
 \label{eq:ns63-relative-control}
\end{equation}
Nevertheless the absolute squared error for this fixed coefficient
diverges:
\begin{equation}
 \norm{\tau_r-c_0\sigma_{1,r}}^2
      \sim\frac{D_0^2}{8\sqrt\pi}\frac{4^r}{r^{3/2}}
      \longrightarrow\infty.
 \label{eq:ns63-absolute-control}
\end{equation}
Moreover $(J/2)^rf\to0$ in $H$ for every fixed $f\in H$.
Neither this normalized convergence nor
\eqref{eq:ns63-relative-control} tests RH.
\end{proposition}
\begin{proof}
The alternating-series expression for zeta shows $z_0<0$: its
numerator $\sum_{n\geq1}(-1)^{n-1}n^{-1/2}$ is positive and its
denominator $1-\sqrt2$ is negative. Thus $c_0$ is well-defined.
Mellin Plancherel writes the relative squared error as the average of
\[
 H(t)=\left|1-\frac{\zeta(1/2+it)}{z_0}\right|^2
\]
under the probability density proportional to
$(1/4+t^2)^{-(r+1)}dt$. The exact second and fourth moments of this
density are
\[
 \mathbb E_r(t^2)=\frac1{4(2r-1)},\qquad
 \mathbb E_r(t^4)=\frac3{16(2r-1)(2r-3)}
\]
for $r\geq1$ and $r\geq2$, respectively. They follow directly
from the beta integral. Near zero, conjugation symmetry and Taylor
expansion give $H(t)=D_0^2t^2+O(t^4)$.
Away from any fixed neighborhood of zero the normalized integral is
exponentially small in $r$. To justify this without an unproved zeta
estimate, use
\[
 \left|\frac{\zeta(1/2+it)}{1/2+it}\right|
 \leq\int_0^\infty\rho_1(x)x^{-1/2}\,dx<4,
\]
or the integrability of the original Mellin square; either controls
the entire exterior integral. The displayed moments then prove
\eqref{eq:ns63-relative-control}. Differentiating the zeta functional
equation at $1/2$, and using
$\psi(1/2)=-\gamma-2\log2$ \cite[Eq. (5.4.13)]{DLMFGammaBounds},
gives the stated value of $D_0$.
Finally $T_r=\binom{2r}{r}\sim4^r/\sqrt{\pi r}$ proves
\eqref{eq:ns63-absolute-control}.

For the last assertion the Mellin multiplier of $(J/2)^r$ is
$(2s)^{-r}$. Its modulus is at most one on $\Re s=1/2$ and tends
to zero for every $t\ne0$. Dominated convergence in the Mellin
$L^2$ norm applies to every fixed $f$. The exceptional point $t=0$
has measure zero.
\end{proof}

\paragraph{Finite control and its full tail.}
The companion Arb calculation integrates the analytic product
\[
 \frac{(1-\zeta(s)/z_0)(1-\zeta(1-s)/z_0)}
             {\pi T_r\{s(1-s)\}^{r+1}},\qquad s=\tfrac12+it,
\]
over $0\leq t\leq A$. It uses reflection $1-s$, not complex
conjugation of the integration variable, so the quadrature callback
is meromorphic. The preceding elementary bound supplies the complete
positive tail bound
\begin{equation}
 \frac1{\pi T_r}\left\{
       \frac{2A^{-2r-1}}{2r+1}
       +\frac{32A^{-2r+1}}{z_0^2(2r-1)}\right\},\qquad r\geq1.
 \label{eq:ns63-quadrature-tail}
\end{equation}
The fixed-order interval results illustrate the control, not the
asymptotic proof or any RH assertion.

The useful research problem is therefore the fixed-order cofinal
correlation estimate in~\eqref{eq:ns61-open-correlation}. Further
smoothing cannot be credited as a convergence result without retaining
one fixed norm and its target. The unsolved arithmetic estimate remains
the one for the optimized coefficients as $N$ increases.

\section{An explicit geometric separator for the positive coefficient cone (NS-64)}
\label{sec:ns64-cone-separator}

The coefficient-cone exclusion in Proposition~\ref{prop:ns62-positive-cone}
has a small, explicit Hilbert-space witness. It quantifies that particular
failure without imposing any restriction on signed NB approximation.
Let $I=(1/6,\infty)$, put $p_1=\sigma_1|_I$ and
$p_n=\epsilon_n|_I$ for $2\leq n\leq6$, and write
\[
 G_I=(\langle p_m,p_n\rangle_{L^2(I)})_{m,n=1}^6,\quad
 \beta=(0,2,3/2,2/3,5/6,-1/5)^T,\quad
 d=G_I^{-1}e_6.
\]
The matrix $G_I$ is positive definite by the nested-support independence
proved in NS-62. In particular $d_6>0$.
The exact local target identity is $\tau|_I=\sum\beta_np_n$.

\begin{proposition}[A quantitative separation, confined to one coefficient cone]
\label{prop:ns64-cone-distance}
Define $\psi=\sum_{n=1}^6d_np_n$ on $I$ and extend it by zero outside
$I$. Then
\[
 \langle\psi,\sigma_1\rangle=0,\qquad
 \langle\psi,\epsilon_n\rangle=\mathbf1_{\{n=6\}}\quad(n\geq2),
 \qquad \langle\psi,\tau\rangle=-\frac15,
 \qquad \norm\psi^2=d_6.
\]
Consequently, if $\mathcal C_+$ denotes the closed cone generated by
$\sigma_1,\epsilon_2,\epsilon_3,\ldots$ with nonnegative real
coefficients, then
\begin{equation}
 \operatorname{dist}_H(\tau,\mathcal C_+)
 \geq\frac1{5\sqrt{d_6}}>0.00228323.
 \label{eq:ns64-positive-gap}
\end{equation}
The exact distance from $\tau|_I$ to the six-generator local cone is
$1/(5\sqrt{d_6})$; the interval replay encloses it between
$0.00228323766118601628$ and $0.00228323766118601629$.
The local equality is not asserted for the full-space distance.
\end{proposition}
\begin{proof}
The defining equation $G_Id=e_6$ gives the first six pairings and
$\norm\psi^2=d^TG_Id=d_6$. Every $\epsilon_n$, $n\geq7$, vanishes
on $I$. The target pairing follows from $\beta_6=-1/5$.
Thus $\langle\psi,f-\tau\rangle\geq1/5$ for every $f\in\mathcal C_+$;
continuity and Cauchy--Schwarz give the first inequality.

For the local equality, put
\[
 c=\beta+\frac{d}{5d_6},\qquad w=\frac{\psi}{5d_6}.
\]
Here $c_6=0$ exactly. The complete interval calculation below certifies
$c_1,\ldots,c_5>0$. Hence $\tau|_I+w=\sum c_np_n$ belongs to the
local cone. Its distance from $\tau|_I$ is $\norm w=1/(5\sqrt{d_6})$,
attaining the dual lower bound. Only the signs of these five finite
coefficients and the displayed decimal enclosure use certified
computation; the separating identity and formula for its bound are exact.
\end{proof}

\paragraph{Complete finite calculation, with no omitted tail.}
Under $t=1/x$, the domain $I$ becomes $0<t<6$ and its measure is
$dt/t^2$. Each $\sigma_n$ becomes $A(t/n)$. On the unit cell
$j<t<j+1$, $k=\lfloor j/n\rfloor$ and
\[
 A(t/n)=t/n-k\log t+k\log n+\log(k!).
\]
On $0<t<1$ the product $A(t/m)A(t/n)/t^2$ is exactly $1/(mn)$;
this cell includes the entire original $x>1$ tail. On the remaining
five cells the products have elementary primitives, retained in the
Maxima and Arb scripts. The triangular difference transformation gives
$G_I$ with every cross term present.

Arb at 256 and 384 bits, using certified preconditioned solves, gives
\[
 c\approx(0.00005096281,\ 1.99987424,\ 1.49456836,\
            0.68321323,\ 0.76905516,\ 0)^T,
\]
and
\[
 \frac1{25d_6}=0.000005213174217458189678\ldots.
\]
The figures in the vector are illustrative rounded displays of the
archived enclosures. All five positivity gates are strict, including
the small first coefficient. The two-precision comparison, complete
Gram and coefficient enclosures, and 39 exact Maxima checks are in
\texttt{evidence/v159/ns64/}. This is a separation of a specified
positive coefficient cone, not negative Weil energy, a lower bound for
the signed best-approximation error, or an obstruction to RH.

\section{Ordinary averaging does not remove the canonical obstruction (NS-65)}
\label{sec:ns65-cesaro}

The preceding results leave different coefficient rules open. One simple
repair of the sharp canonical rule can be tested exactly: ordinary
Ces\`aro averaging over its index. The conclusion below concerns a fixed
number of these averages. Logarithmic tapers, a growing number of averages,
and optimized coefficients are outside its scope.

For a sequence $(f_N)$ of vectors or scalars define
\[
 (Cf)_N=\frac1N\sum_{n=1}^Nf_n,\qquad
 B_N^{[q]}=(C^qB)_N,\qquad
 H_N^{[q]}=(C^q a)_N,\quad a_n=n\eta_n,
\]
where $B_n$ and $\eta_n$ are the canonical objects of NS-62 and
$q\geq0$ is a fixed integer. The averaging weights are nonnegative,
sum to one, and only involve indices $n\leq N$.

\begin{proposition}[A fixed number of Ces\`aro averages leaves a critical-zero pole]
\label{prop:ns65-cesaro-obstruction}
For every fixed integer $q\geq0$ and every sequence of nonnegative
integers $r_N$,
\[
 \norm{J^{r_N}(\chi+B_N^{[q]})}_H\not\longrightarrow0
 \quad(N\to\infty\text{ through all integers}).
\]
More precisely, for every $N$ and every $r\geq0$,
\begin{equation}
 \norm{J^r(\chi+B_N^{[q]})}_H
 \geq\frac{|H_N^{[q]}-k_{N,r}|}{\sqrt N}
 \geq\frac{|H_N^{[q]}-2|}{\sqrt N}
             -\frac{4\pi^2}{3^{r+2}N^{5/2}},
 \label{eq:ns65-averaged-witness}
\end{equation}
and $H_N^{[q]}\ne o(\sqrt N)$. This is full-sequence nonconvergence;
no selected-subsequence exclusion or assertion about signed optimal
approximation follows.
\end{proposition}
\begin{proof}
For $0<x<1/N$, every index $n\leq N$ obeys the same small-cell
identity from NS-62. Linearity and normalization of the weights give
\[
 \mathcal V(\chi+B_N^{[q]})(x)=k(x)-H_N^{[q]}.
\]
Apply the same unit Laguerre test $\psi_{N,r}$ and its complete
adjoint identity. This proves~\eqref{eq:ns65-averaged-witness},
including its remainder uniform in $r$.

It remains to prove the scalar obstruction. Put
\[
 \eta(x)=\int_1^x\frac{M(t)}{t^2}\,dt,\qquad
 A_0(x)=x\eta(x),\qquad
 A_{q+1}(x)=\frac1x\int_1^x A_q(u)\,du.
\]
For $\Re s>1$, absolute convergence and Fubini give
\begin{equation}
 \int_1^\infty A_q(x)x^{-s-1}\,dx
   =\frac1{(s+1)^q(s-1)s\zeta(s)}.
 \label{eq:ns65-cesaro-mellin}
\end{equation}
For $q=0$ this is the Mellin identity for $\eta$ with its variable
shifted by one. Each continuous averaging contributes exactly the
factor $(s+1)^{-1}$; there is no discarded endpoint term.

We now connect this continuous identity to the discrete averages.
Directly $A_0(x)=O(x\log(2x))$ and, almost everywhere,
\[
 A_0'(x)=\eta(x)+M(x)/x
         =\sum_{n\leq x}\mu(n)/n=O(\log(2x)).
\]
Inductively $A_q(x)=O_q(x\log(2x))$ and
$A_{q+1}'(x)=(A_q(x)-A_{q+1}(x))/x=O_q(\log(2x))$.
The functions are locally absolutely continuous. Comparing each
unit-cell integral with its endpoint value therefore gives
\[
 \frac1N\sum_{n\leq N}A_q(n)-\frac1N\int_1^N A_q(x)\,dx
       =O_q(\log(2N)).
\]
Since $A_0(N)=a_N$, induction yields
\begin{equation}
 H_N^{[q]}=A_q(N)+O_q(\log(2N)).
 \label{eq:ns65-discrete-continuous}
\end{equation}
Thus $H_N^{[q]}=o(\sqrt N)$ would imply $A_q(x)=o(\sqrt x)$
also between integers, by the displayed derivative bounds.

Under that hypothetical little-$o$ estimate the Mellin integral in
\eqref{eq:ns65-cesaro-mellin} is analytic for $\Re s>1/2$.
At a known critical-line zero $\rho=1/2+i\gamma_*$, its usual
Abelian boundary estimate gives
\[
 \lim_{\sigma\downarrow1/2}(\sigma-1/2)
       \int_1^\infty A_q(x)x^{-\sigma-i\gamma_*-1}\,dx=0.
\]
For example, split the integral at a fixed large $X$, use
$|A_q(x)|\leq\varepsilon\sqrt x$ on the tail, then let
$\sigma\downarrow1/2$ and $\varepsilon\downarrow0$.
But the meromorphic continuation on the right side of
\eqref{eq:ns65-cesaro-mellin} has a genuine pole at $\rho$:
none of $\rho+1$, $\rho-1$ or $\rho$ is zero, and its numerator
is one. A pole of any positive order contradicts the boundary limit.
Therefore $\limsup_N|H_N^{[q]}|/\sqrt N>0$.
Together with~\eqref{eq:ns65-averaged-witness}, whose remainder tends
to zero uniformly in $r\geq0$, this proves the stated nonconvergence.
\end{proof}

\paragraph{Interpretation and stop condition.}
Ordinary averaging attenuates a critical-zero contribution by
$(\rho+1)^{-q}$, but a fixed $q$ cannot annihilate it. This is an
analytic obstruction for the specified averaged sharp rule, not a
numerical observation or a priority claim about summability methods.
The general summability viewpoint is already part of
\cite[Section 3]{BaezDuarteStrong}. The useful open task remains an
independent estimate for the optimized signed residual, or a genuinely
different coefficient rule with a complete norm bound. No such lower
correlation or convergence estimate is supplied here.

\section{A weaker canonical target: growth rather than convergence (NS-66)}
\label{sec:ns66-canonical-growth}

Strong convergence of the full canonical sequence is excluded, but
that does not exclude slow growth or bounded subsequences of its norms.
One Mellin integration makes a complete Abel-summation estimate
available. The resulting growth exponent is an exact reformulation
of the zero-location question, not an independently established RH bound.

For a fixed integer $r\geq1$ define
\[
 D_{N,r}=\norm{J^r(\chi+B_N)}_H,\qquad
 \beta_* =\sup\{\Re z:\zeta(z)=0,\ 0<\Re z<1\}.
\]
Known critical-line zeros give $1/2\leq\beta_*\leq1$.

\begin{proposition}[Complete growth bounds for the smoothed canonical error]
\label{prop:ns66-growth-bounds}
If $z=\beta+i\gamma_*$ is a zero with $1/2<\beta<1$, then
\begin{equation}
 D_{N,r}\geq\frac{\sqrt{2\beta-1}}{|z|^{r+1}}N^{\beta-1/2}
 \quad\text{for every }N\geq1.
 \label{eq:ns66-zero-growth}
\end{equation}
Conversely, suppose for some $1/2\leq\theta<1$ and constant $C$
that $|M(u)|\leq Cu^\theta$ for all $u\geq1$. Then
\begin{equation}
 D_{N,1}\leq\sqrt2+
  \frac{C\norm{\sigma_1}}{1-\theta}N^{\theta-1/2}
  +C\sqrt\kappa\int_1^N u^{\theta-3/2}\,du,
 \qquad D_{N,r}\leq2^{r-1}D_{N,1}.
 \label{eq:ns66-mertens-upper}
\end{equation}
The integral is $(N^{\theta-1/2}-1)/(\theta-1/2)$ for
$\theta>1/2$, and $\log N$ for $\theta=1/2$.
The Mertens bound in this implication is a hypothesis, not a result
proved here.
\end{proposition}
\begin{proof}
Exact exterior interpolation gives $\chi+B_N=0$ on $x>1/N$.
Since $Jf(x)$ only uses arguments at least $x$, $J^r(\chi+B_N)$
is supported in $(0,1/N]$. Its Mellin value at $z$ is
$z^{-(r+1)}$: the zeta factor annihilates every basis contribution.
Cauchy--Schwarz on that support gives
\[
 |z|^{-(r+1)}\leq D_{N,r}
 \left(\int_0^{1/N}x^{2\beta-2}\,dx\right)^{1/2}
 =D_{N,r}\frac{N^{1/2-\beta}}{\sqrt{2\beta-1}},
\]
proving~\eqref{eq:ns66-zero-growth}.

Under the stated Mertens hypothesis, partial summation makes
$\sum\mu(n)/n$ convergent and identifies its value as
$\lim_{s\downarrow1}1/\zeta(s)=0$. Hence
\[
 \eta_N=-\int_N^\infty M(u)u^{-2}\,du,\qquad
 |\eta_N|\leq\frac{C}{1-\theta}N^{\theta-1}.
\]
The strong derivative $\partial_u\sigma_u=-\rho_u/u$ gives the
complete Bochner summation-by-parts identity
\begin{equation}
 JB_N=-N\eta_N\sigma_N+
           \int_1^N\frac{M(u)}u\rho_u\,du.
 \label{eq:ns66-bochner-abel}
\end{equation}
There is no lower-endpoint contribution, since $M(1^-)=0$.
Use $\norm{\sigma_N}=N^{-1/2}\norm{\sigma_1}$,
$\norm{\rho_u}=\sqrt{\kappa/u}$ and $\norm{J\chi}=\sqrt2$.
The triangle inequality yields~\eqref{eq:ns66-mertens-upper},
including all coefficient and tail-correction terms. Boundedness
of $J^{r-1}$ gives the final assertion.
\end{proof}

\begin{corollary}[The canonical norm-growth exponent records the rightmost zeros]
\label{cor:ns66-growth-exponent}
For every fixed integer $r\geq1$ the following limit exists and equals
\begin{equation}
 \lim_{N\to\infty}\frac{\log(1+D_{N,r})}{\log N}
                  =\beta_*-\frac12.
 \label{eq:ns66-exponent}
\end{equation}
In particular, RH is equivalent to the subpolynomial bound
\[
 \text{for every }\varepsilon>0,\qquad
 D_{N,r}=O_{r,\varepsilon}(N^\varepsilon).
\]
Even a bounded cofinal subsequence of $D_{N,r}$ would suffice to
prove RH. No such bounded subsequence or unconditional subpolynomial
bound is supplied, and boundedness is not claimed necessary under RH.
\end{corollary}
\begin{proof}
For each zero to the right of the line,
\eqref{eq:ns66-zero-growth} bounds the lower limit by
$\beta-1/2$. Taking the supremum gives $\beta_*-1/2$; when
$\beta_*=1/2$, the same lower limit is at least zero trivially.

If $\beta_*<1$, choose $\beta_*<\alpha<\theta<1$.
The classical zero-free-half-plane implication, stated precisely
in~\cite[Lemma 2.1]{BaezDuarteStrong}, gives convergence of
$\sum\mu(n)n^{-\theta}$. Partial summation of its bounded partial
sums gives $M(x)=O_\theta(x^\theta)$.
Proposition~\ref{prop:ns66-growth-bounds} then bounds the upper
limit by $\theta-1/2$. Let $\theta\downarrow\beta_*$.
This includes RH by choosing $\theta>1/2$ arbitrarily close to the
line, without claiming a bound at the endpoint $\theta=1/2$.

If $\beta_*=1$, the direct bounds $|s_N|\leq1+\log N$ and
$\norm{\sigma_n}=\norm{\sigma_1}/\sqrt n$ give
$D_{N,1}=O(\sqrt N\log(2N))$, so the upper limit is at most
$1/2$. This matches the preceding lower bound.
Thus~\eqref{eq:ns66-exponent} holds in every case.
The functional equation identifies $\beta_*=1/2$ with RH.
Finally~\eqref{eq:ns66-zero-growth} diverges along every cofinal
subsequence if any zero lies to the right of the line, proving the
bounded-subsequence implication.
\end{proof}

\paragraph{What the weaker target changes.}
For this explicit coefficient rule, norm convergence is false, while
subpolynomial norm growth is equivalent to RH. These are compatible
statements. The exact growth reformulation does not itself improve
our knowledge of $\beta_*$ or supply the missing Mertens estimate.
It gives a second precise research target: a direct subpolynomial
bound for the complete smoothed canonical norm, or a bounded cofinal
subsequence. Increasing the number of Mellin integrations changes
constants in~\eqref{eq:ns66-zero-growth}; the corollary is stated
only for one fixed order. No inference from a changing normalization
or finite table is used.
